\documentclass[10pt,a4paper]{amsart}
\usepackage[margin=19mm]{geometry}
\usepackage{amsmath,amssymb,mathtools}
\usepackage[scr=rsfs,cal=euler]{mathalfa}
\usepackage{xfrac}
\usepackage[unicode,hidelinks,bookmarks=false]{hyperref}
\newcommand{\ie}{i.e.\ }
\newcommand{\cf}{cf.\ }
\newcommand{\resp}{respectively\ }
\newcommand{\Subsection}{\subsection}
\providecommand{\nobreakdash}{\nobreak\hbox{-}}
\setlength{\emergencystretch}{2em}
\multlinegap0pt
\usepackage{amssymb}
\usepackage{mathtools}
\usepackage{tikz}
\newtheorem{proposition}{Proposition}
\title{Exact Uniform L1 Spacing for Solow--Polasky Diversity on Lines and Ordered Pareto Fronts}
\author{Michael T. M. Emmerich, Mahboubeh Nezhadmoghaddam,, Jes\'us Guillermo Falc\'on Cardona}
\date{Source: arXiv:2605.21922v1 (CC BY 4.0)}
\begin{document}
\maketitle

\noindent\emph{Source and licence.} Exact Uniform L1 Spacing for Solow--Polasky Diversity on Lines and Ordered Pareto Fronts. Version arXiv:2605.21922v1 (CC BY 4.0), distributed by arXiv under the Creative Commons Attribution 4.0 International licence, http://creativecommons.org/licenses/by/4.0/, as recorded in the arXiv licence metadata for this exact version and shown on the abstract page https://arxiv.org/abs/2605.21922v1. Full licence notice: \texttt{licenses/arxiv-2605.21922-CC-BY-4.0.txt}.\par\smallskip

\noindent\textit{Editorial scope.} Counted unit: the article's proposition prop:converse-exponential-kernel (source0.tex lines 265--288). The article's complete original proof of it is delivered with it (source0.tex lines 290--347, one \texttt{proof} environment of 58 lines). No mathematics is written, repaired or re-derived: the statement and the proof are the article's own lines, in the article's own notation, and no retained line is substituted or re-flowed.  No further source-local context is needed: every cross-reference of every retained line resolves inside the two delivered spans.  Nothing was omitted from any retained span, so the package carries no omission record.  No retained line contains a citation, so the wrapper carries no bibliography and no bibliography entry.  No retained line contains a figure environment, an image-inclusion command or drawing code, so no image is delivered and the package declares no asset dependency.  Editorial formatting only, in addition: an AMS \texttt{amsart} preamble that restates the article's own declarations and macros used by the retained lines, namely the environments \texttt{proposition} on separate counters, together with the standard packages those lines need.  The retained environments print in this document as Proposition 1 in order of appearance, whereas the article prints the same items with the numbering of its own full document.  The complete native source of the article as distributed in the arXiv e-print for this exact version is bundled unchanged as \texttt{sources/201099/source0.tex}.
\begin{proposition}[Adjacent-gap additivity forces the exponential family]
\label{prop:converse-exponential-kernel}
Let $K:[0,\infty)\to(0,1]$ be continuous, non-increasing, normalized by $K(0)=1$, and satisfy $K(t)<1$ for every $t>0$.  For a finite ordered set $S=\{x_1<\cdots<x_k\}\subset\mathbb{R}$, define
\[
  (Z_K(S))_{ij}=K(|x_i-x_j|),
  \qquad
  D_K(S)=\mathbf{1}^\top Z_K(S)^{-1}\mathbf{1},
\]
whenever $Z_K(S)$ is nonsingular, and put $E_K(S)=D_K(S)-1$.  Suppose that this excess diversity is additive in adjacent gaps for every three-point set:
\begin{equation}
\label{eq:three-point-excess-additivity}
  E_K(\{0,a,a+b\})
  =E_K(\{0,a\})+E_K(\{0,b\})
  \qquad (a,b>0).
\end{equation}
Then $K$ belongs to the exponential family.  More precisely,
\[
  K(t)=e^{-\beta t}\qquad(t\geq0),
\]
where the parameter is uniquely fixed by the one-step value,
\[
  \beta=-\log K(1)>0.
\]
\end{proposition}

\begin{proof}
For two points at distance $t$, writing $r=K(t)$ gives
\[
  D_K(\{0,t\})
  =\mathbf{1}^\top
  \begin{pmatrix}1&r\\ r&1\end{pmatrix}^{-1}
  \mathbf{1}
  =\frac{2}{1+r},
\]
and hence
\[
  E_K(\{0,t\})=\frac{1-r}{1+r}.
\]
Fix $a,b>0$ and write
\[
  u=K(a),\qquad v=K(b),\qquad w=K(a+b).
\]
For the three-point set $\{0,a,a+b\}$, the similarity matrix is
\[
  Z_K=
  \begin{pmatrix}
  1 & u & w\\
  u & 1 & v\\
  w & v & 1
  \end{pmatrix}.
\]
Substituting this matrix and the two-point formula into~\eqref{eq:three-point-excess-additivity}, and simplifying, gives
\[
  (uv-w)\bigl(u^2-uvw-uv+v^2+w-1\bigr)=0.
\]
Thus either $w=uv$, or
\[
  w=\frac{1+uv-u^2-v^2}{1-uv}.
\]
The second alternative is impossible under the non-increasing assumption.  Interchanging $a$ and $b$ if necessary, assume $u\geq v$.  Since $a,b>0$, the assumptions give $0<v\leq u<1$.  Since $a+b>b$ and $K$ is non-increasing, we also have $w\leq v$.  However, the second alternative gives
\[
  w-v
  =\frac{(1-u)(u+1-v-v^2)}{1-uv}>0,
\]
because $u\geq v$ implies $u+1-v-v^2\geq1-v^2>0$.  This contradicts $w\leq v$.  Therefore only the first alternative remains, and
\[
  K(a+b)=K(a)K(b)
  \qquad(a,b>0).
\]
Together with $K(0)=1$, this identity holds on $[0,\infty)$.

Since $K$ is positive, $h(t)=\log K(t)$ is continuous and satisfies
\[
  h(a+b)=h(a)+h(b).
\]
The continuous Cauchy equation on $[0,\infty)$ gives $h(t)=t h(1)$.  Hence
\[
  K(t)=e^{t\log K(1)}=e^{-\beta t},
  \qquad
  \beta=-\log K(1).
\]
Because $K(1)<1$, we have $\beta>0$, and the displayed formula shows that this value of $\beta$ is unique for the given kernel.
\end{proof}

\end{document}
